\chapter{绪\quad 论}

\section{研究背景与意义}

\subsection{研究背景}

近年来，智能可穿戴设备进入快速普及阶段。以 Ray-Ban Meta 智能眼镜\textsuperscript{[1]}和 Apple Vision Pro 头戴式空间计算设备\textsuperscript{[2]}为代表的新一代产品，已经具备高分辨率摄像头、片上神经网络推理引擎和持续联网能力，可以在佩戴者几乎无任何外显动作的情况下完成对周围环境的实时图像采集与语义理解。这类设备获取屏幕信息的方式主要有两种：一是抽帧分析（Snapshots），即每隔固定时间截取一张高分辨率图片上传云端进行光学字符识别（Optical Character Recognition, OCR）或图像特征匹配；二是流式识别（Streaming），即以较低分辨率连续处理视频流，用于实时翻译、场景理解等任务。

在办公室、实验室、银行柜台等场景中，显示器上的源代码、财务报表、账号口令、个人隐私数据等敏感信息，在上述两种采集方式面前几乎是完全暴露的：攻击者只需佩戴一副外观与普通眼镜无异的智能眼镜，从屏幕前走过即可完成信息窃取，事后亦难以取证。传统的信息安全体系着重保护数据在存储与传输环节的机密性，而"屏幕显示"这一信息呈现的最后环节，恰恰是目前防护最薄弱的环节。

现有的显示器防偷拍技术主要分为三类，但均存在明显不足。第一类是物理遮蔽方案，如防窥膜、光学百叶窗滤镜等，依赖额外的硬件屏蔽结构，成本高、安装不便，并且会压缩使用者自身的可视角度、降低色彩还原度，难以在办公环境中大规模部署。第二类是简单高频闪烁方案，通过全屏高频亮灭交替干扰相机采样，但该方法调制深度大、易导致人眼疲劳，且现代识别算法可以通过多帧时域平均或自适应曝光轻松破解。第三类是基于光源干扰的主动防御方案，如在屏幕周围部署红外 LED 阵列，需要额外硬件、功耗较高，且对不同型号相机的干扰效果差异很大，缺乏普适性。

\subsection{研究意义}

针对上述问题，本课题提出并实现一种纯软件的显示器隐私保护方案：利用人眼视觉暂留（Persistence of Vision, POV）效应与相机高速快门采样之间的根本性差异，将每一帧完整图像在时间维度上拆分为多个互补的随机点阵子帧并以高刷新率轮流显示，使人眼在积分窗口内感知到完整画面，而相机在任一快门时刻只能捕获碎片化残片。

本课题的研究意义体现在三个方面。其一，方案无需对显示器进行任何硬件改造，仅依赖软件渲染链路即可部署，克服了物理遮蔽方案成本高、影响观看体验的缺点；其二，方案将密码学安全伪随机掩模与对抗样本噪声相结合，同时面向"人眼无感"与"机器不可识别"两个对立目标进行联合设计，对抽帧分析与流式识别两类真实威胁均有量化验证的防御效果；其三，本课题以概念验证（Proof of Concept, PoC）系统的形式完整实现了掩模生成、噪声注入、时序控制、渲染合成四个核心模块，并通过攻防实验诚实地刻画了线性时域方案的能力边界，为后续驱动层产品化与专利申请提供了可靠的实验依据。

\section{国内外研究现状}

\subsection{国外研究现状}

屏幕与相机之间的攻防问题在学术界已有持续研究。Zhang 等人\textsuperscript{[3]}在 MobiCom 会议上提出了 Kaleido 系统，通过分析人眼视觉与相机成像在时域响应上的差异，对视频帧进行时域调制，实现"可观看但不可录制"的版权保护效果，验证了利用视觉暂留差异进行显示端防护的可行性，但其目标是视频防盗录，调制方式以帧间颜色扰动为主，未针对静态文本的 OCR 威胁设计像素级掩模。Zhu 等人\textsuperscript{[4]}提出 LiShield 系统，利用智能 LED 光源的高频调制在拍摄图像中引入条带失真以保护物理空间隐私，属于光源干扰类主动防御，需要部署专用光源硬件。Fang 等人\textsuperscript{[5]}研究了抗屏摄水印（screen-shooting resilient watermarking），关注的是拍屏行为发生后的溯源取证，而非阻止信息被识别。

在对抗机器识别方面，Goodfellow 等人\textsuperscript{[6]}提出快速梯度符号法（Fast Gradient Sign Method, FGSM），证明在图像上叠加人眼难以察觉的有界扰动即可使深度神经网络产生错误输出；Madry 等人\textsuperscript{[7]}进一步提出投影梯度下降法（Projected Gradient Descent, PGD），成为目前公认的一阶对抗攻击强基线；Carlini 与 Wagner\textsuperscript{[8]}的工作则系统评估了神经网络在对抗扰动下的脆弱性。Akhtar 等人\textsuperscript{[9]}对计算机视觉中的对抗攻击进行了系统综述。这些工作为本课题"以对抗噪声最大化单帧识别错误率"的设计提供了理论与算法基础。

\subsection{国内研究现状}

国内在显示安全方向的研究与产业实践主要集中在防窥膜等光学材料、屏幕水印溯源以及涉密场所的管理性措施上。市售防窥膜通过微百叶窗结构限制可视角度，对正面拍摄无能为力；企业级终端普遍采用的屏幕水印方案在信息被拍摄识别后只能用于事后追责。近年来随着对抗样本研究的兴起，国内学术界在对抗扰动生成、OCR 鲁棒性评估等方向产出了大量成果，但将对抗噪声与高刷新率时域掩模相结合、面向智能眼镜实时拍摄威胁的软件级显示防护方案，目前尚缺乏公开的系统性实现与量化评测。

\subsection{发展趋势}

综合国内外研究现状可以看出三个趋势：第一，攻击端能力持续增强，智能眼镜的端侧算力与云端识别服务使"随手拍摄即识别"成为常态，被动的物理遮蔽难以应对；第二，防御端从"遮挡可见光路径"向"调制显示时序"演进，利用人眼与相机感知差异的时域方案成为研究热点；第三，单一手段难以同时对抗多种采集策略，掩模、噪声、时序、亮度补偿等多机制联合设计成为必然方向。本课题正是顺应上述趋势，将密码学随机掩模、对抗噪声与高刷新率时序控制融合为一个完整系统。

\section{研究目标与主要内容}

\subsection{研究目标}

本课题的总体目标是：设计并实现一个基于视觉暂留与时间分片掩模的显示器隐私保护概念验证系统，使人眼能够正常、无感知地阅读屏幕完整信息，而采用抽帧分析或流式识别策略的相机设备在任一瞬时采样中均无法获取可被 OCR 或目标检测算法识别的图像信息。具体需同时满足四项约束：

（1）视觉完整性约束：人眼通过约 50ms 时间窗口的视觉暂留积分能够感知到与原始图像无差别的完整画面；

（2）瞬时不可识别约束：任意单帧仅包含原始图像的随机碎片化残片，识别准确率降至接近零；

（3）抗叠加约束：攻击者在不足一个完整周期内进行多帧叠加，或采用卷帘快门混合采样，均无法获得有效信息；

（4）抗长曝光约束：攻击者延长曝光时间试图积分还原时，被插入的反色帧或黑帧有效对抗。

\subsection{研究内容}

围绕上述目标，本课题的主要研究内容包括：

（1）随机点阵掩模生成算法。基于 ChaCha20 密码学安全伪随机数生成器构造随机索引矩阵，将图像拆分为满足完备性约束 $\sum I_k = I$ 与互斥性约束的 $n$ 个子帧，并通过 HKDF 密钥派生保证相邻周期统计独立、通过卡方检验保证掩模分布均匀。

（2）时域互补对抗噪声注入算法。基于 FGSM/PGD 生成针对 OCR 模型的有界对抗扰动，将其分解为满足 $\sum N_k = 0$ 的互补子噪声，使噪声在人眼积分后完全抵消、在单帧内最大化识别错误率，并支持多目标模型轮换与在线策略更新。

（3）时序控制与渲染合成。实现 Fisher-Yates 伪随机子帧置换、VBlank 时序调度、反色帧与黑帧插入、亮度补偿、HDR 色域处理、多显示器同步以及 GPU 着色器渲染链路。

（4）攻防量化评测体系。构建全局快门、卷帘快门、时域平均、长曝光等相机攻击模拟器，使用 Tesseract、EasyOCR、PaddleOCR 三种 OCR 引擎与 YOLOv8 目标检测模型在多样本语料上量化防御效果，并以闪烁感知指数、CIEDE2000 色差等指标验证人眼无感性。

\subsection{技术路线}

本课题采用"理论建模—模块实现—单元验证—攻防实验—诚实评估"的技术路线，如图1-1所示。首先对视觉暂留积分、掩模完备性、噪声互补性进行数学建模；然后以 Python 实现核心模块并配套 118 项单元测试保证算法正确性；继而构建相机攻击模拟器开展端到端攻防实验；最后将实测结果与设计目标逐项对照，明确方案的有效边界与原理性局限。

\begin{figure}[htbp]
  \centering
  \begin{tikzpicture}[
      node distance=4mm and 6mm,
      box/.style={rectangle, draw, rounded corners=1pt, minimum height=9mm,
                  minimum width=24mm, align=center, font=\zihao{5}},
      arr/.style={-{Stealth[length=2mm]}, thick}]
    \node[box] (a) {理论建模\\掩模/噪声/时序};
    \node[box, right=of a] (b) {模块实现\\四核心模块};
    \node[box, right=of b] (c) {单元验证\\118 项测试};
    \node[box, right=of c] (d) {攻防实验\\OCR/检测};
    \node[box, right=of d] (e) {诚实评估\\边界与局限};
    \draw[arr] (a) -- (b);
    \draw[arr] (b) -- (c);
    \draw[arr] (c) -- (d);
    \draw[arr] (d) -- (e);
  \end{tikzpicture}
  \caption{技术路线}
\end{figure}

\section{论文组织结构}

本文共分为五章。第一章绪论，介绍课题背景意义、国内外研究现状与研究内容；第二章介绍视觉暂留、密码学伪随机数、对抗样本与 OCR 识别等相关理论与关键技术；第三章进行威胁模型与需求分析，给出系统总体架构与核心数学模型；第四章详细阐述掩模生成器、噪声注入器、渲染驱动器、时序控制器等核心模块的设计与实现；第五章给出视觉无感性评估、攻防实验、消融实验与性能实测结果，并分析方案的局限性；最后是结论与展望。

\section{本章小结}

本章阐述了智能眼镜普及背景下显示器信息泄露的安全威胁，分析了物理遮蔽、高频闪烁、光源干扰三类现有方案的不足，综述了国内外在显示端时域防护与对抗样本方向的研究现状，明确了本课题"人眼无感、相机不可识别"的研究目标、四项设计约束、主要研究内容与技术路线，最后给出了论文的组织结构。
